% ============================================================
% SPRINT 4: ROLE-BASED WORKFLOWS AND COMMUNITY FEATURES
% ============================================================
\section{Sprint 4: Role-Based Workflows and Community Features}

\subsection*{Introduction}
Sprint 4 extends the platform with features tailored to specific user roles: a structured intern access workflow, seller analytics, asset edit request management, peer-to-peer asset trading, and community groups. This sprint makes the RBAC system operationally meaningful for every user segment.

\subsection{Sprint Backlog}

\begin{sprintbox}{Sprint 4 --- Duration: 3 weeks}
\textbf{Goal:} Deliver role-specific advanced features for interns, sellers, and the broader community.
\end{sprintbox}

\vspace{0.4cm}

\begin{table}[H]
\centering
\caption{Sprint 4 --- User Stories}
\begin{tabularx}{\textwidth}{C{1cm} X C{2cm} C{2cm}}
\toprule
\textbf{ID} & \textbf{User Story} & \textbf{Priority} & \textbf{Points} \\
\midrule
US-24 & As an intern, I want to request access to a paid asset for IGS projects. & High & 5 \\
US-25 & As an admin, I want to approve intern asset requests and generate vouchers. & High & 5 \\
US-26 & As an intern, I want to redeem a voucher code to unlock an asset. & High & 3 \\
US-27 & As a seller, I want to view my sales analytics (views, downloads, revenue). & High & 5 \\
US-28 & As a seller, I want to submit an edit request for an approved asset. & Medium & 5 \\
US-29 & As a user, I want to propose a trade with another user. & Medium & 8 \\
US-30 & As a user, I want to create and join community groups. & Medium & 5 \\
US-31 & As an external user, I want to request the seller role. & High & 3 \\
\midrule
\multicolumn{3}{r}{\textbf{Total Story Points:}} & \textbf{39} \\
\bottomrule
\end{tabularx}
\end{table}

\subsection{Analysis and Design}

\subsubsection{Intern Asset Request Workflow}

\begin{figure}[H]
\centering
\resizebox{\textwidth}{!}{% ============================================================
% Sequence Diagram - Intern Asset Request + Voucher Redemption
% (Visual Paradigm style)
% ============================================================
\begin{tikzpicture}[
  font=\sffamily\scriptsize,
  part/.style={rectangle, draw=darkGray, line width=0.6pt,
               fill=yellow!18, text=black,
               minimum width=2.6cm, minimum height=0.7cm,
               font=\sffamily\footnotesize\bfseries,
               align=center, inner sep=3pt},
  ll/.style={dashed, draw=darkGray!70, line width=0.45pt},
  activ/.style={rectangle, draw=darkGray, line width=0.5pt, fill=isiBlue!12},
  sync/.style ={-{Stealth[length=2.2mm]}, draw=darkGray, line width=0.7pt},
  ret/.style  ={-{Latex[length=2mm]},     draw=darkGray, line width=0.55pt, dashed},
  note/.style={rectangle, draw=orange!70!black, line width=0.5pt,
               fill=yellow!30, font=\scriptsize, align=left,
               inner sep=3pt, rounded corners=1pt},
  phase/.style={rectangle, draw=isiBlue, line width=0.6pt,
                fill=isiBlue!10, font=\sffamily\scriptsize\bfseries\color{isiBlue},
                align=center, inner sep=3pt, rounded corners=2pt},
  lbl/.style={font=\scriptsize, fill=white, inner sep=1pt}
]

  % ---- Participants ----
  \node[part] (intern)  at (0,   0) {Intern};
  \node[part] (api)     at (4.5, 0) {Backend\\API};
  \node[part] (db)      at (9.0, 0) {PostgreSQL};
  \node[part] (admin)   at (13.0,0) {Administrator};
  \node[part] (notif)   at (16.5,0) {Notification\\Service};

  % ---- Lifelines ----
  \foreach \x in {0, 4.5, 9.0, 13.0, 16.5}{
    \draw[ll] (\x,-0.4) -- (\x,-15.5);
  }

  % ---- Activation boxes ----
  \filldraw[activ] (-0.12,-0.6)   rectangle (0.12,-15.0);
  \filldraw[activ] (4.38,-1.0)    rectangle (4.62,-4.0);
  \filldraw[activ] (8.88,-2.0)    rectangle (9.12,-3.5);
  \filldraw[activ] (4.38,-5.5)    rectangle (4.62,-10.5);
  \filldraw[activ] (8.88,-6.5)    rectangle (9.12,-9.5);
  \filldraw[activ] (12.88,-5.2)   rectangle (13.12,-5.6);
  \filldraw[activ] (4.38,-12.0)   rectangle (4.62,-14.8);
  \filldraw[activ] (8.88,-13.0)   rectangle (9.12,-14.4);

  % =====================================================================
  % PHASE 1 — Intern submits request
  % =====================================================================
  \node[phase, anchor=east] at (-0.3,-1.0) {Phase 1\\Submit\\Request};

  \draw[sync] (0.12,-1.0) -- node[lbl] {1: POST /api/intern/requests \{assetId, justification\}} (4.38,-1.0);
  \draw[sync] (4.62,-2.0) -- node[lbl] {2: INSERT intern\_asset\_requests (status=pending)} (8.88,-2.0);
  \draw[ret]  (8.88,-3.0) -- node[lbl] {request id} (4.62,-3.0);
  \draw[sync] (4.62,-3.5) -- node[lbl] {3: notify(admins, "intern\_request")} (16.38,-3.5);
  \draw[ret]  (4.38,-3.9) -- node[lbl] {201 \{ submitted \}} (0.12,-3.9);

  % =====================================================================
  % PHASE 2 — Administrator approves (asynchronous)
  % =====================================================================
  \node[phase, anchor=west] at (16.8,-5.5) {Phase 2\\Admin\\Approves};

  \draw[sync] (12.88,-5.5) -- node[lbl] {4: POST /api/admin/intern-requests/:id/approve} (4.62,-5.5);
  \draw[sync] (4.62,-6.5) -- node[lbl] {5: generate voucher code} (8.88,-6.5);
  \draw[sync] (4.62,-7.2) -- node[lbl] {6: INSERT intern\_vouchers (code)} (8.88,-7.2);
  \draw[sync] (4.62,-8.0) -- node[lbl] {7: UPDATE intern\_asset\_requests SET status=approved} (8.88,-8.0);
  \draw[ret]  (8.88,-8.9) -- node[lbl] {ok} (4.62,-8.9);
  \draw[sync] (4.62,-9.4) -- node[lbl] {8: notify(intern, "approved", voucherCode)} (16.38,-9.4);
  \draw[ret]  (4.38,-10.3) -- node[lbl] {200 \{ voucherCode \}} (12.88,-10.3);

  % =====================================================================
  % PHASE 3 — Intern redeems voucher
  % =====================================================================
  \node[phase, anchor=east] at (-0.3,-12.0) {Phase 3\\Redeem\\Voucher};

  \draw[sync] (0.12,-12.0) -- node[lbl] {9: POST /api/intern/redeem-voucher \{code\}} (4.38,-12.0);
  \draw[sync] (4.62,-13.0) -- node[lbl] {10: SELECT intern\_vouchers WHERE code AND is\_used=false} (8.88,-13.0);
  \draw[ret]  (8.88,-13.8) -- node[lbl] {valid voucher} (4.62,-13.8);
  \draw[sync] (4.62,-14.3) -- node[lbl] {11: INSERT user\_downloads + UPDATE voucher.is\_used} (8.88,-14.3);
  \draw[ret]  (4.38,-14.8) -- node[lbl] {200 \{ asset unlocked \}} (0.12,-14.8);

  % ---- Termination ----
  \foreach \x in {0, 4.5, 9.0, 13.0, 16.5}{
    \draw[line width=0.8pt] (\x-0.15,-15.05) -- (\x+0.15,-15.05);
  }

\end{tikzpicture}
}
\caption{Intern Asset Request Workflow (Sequence Diagram)}
\label{fig:sprint4_intern_flow}
\end{figure}

The intern workflow proceeds as follows: the intern submits a request with a written justification; the admin reviews it and optionally generates a single-use voucher code; the intern redeems the voucher, which unlocks the asset in their personal library. New database tables: \texttt{intern\_asset\_requests} and \texttt{intern\_vouchers}.

\subsubsection{Peer-to-Peer Trading System}

The trading system allows two asset owners to propose an exchange without monetary transaction. User A offers asset X in exchange for asset Y owned by User B. If User B accepts, both assets are atomically added to each other's libraries and the \texttt{trades} record is marked \texttt{completed}.

\subsubsection{Seller Analytics}

The analytics dashboard aggregates: total revenue, total orders, total downloads, active asset count, and a per-asset breakdown (revenue, download count, average rating) — all computed from a single aggregate SQL query joining \texttt{assets}, \texttt{order\_items}, and \texttt{user\_downloads} tables.

\begin{figure}[H]
\centering
\resizebox{\textwidth}{!}{% ER Diagram -- Sprint 4: Intern Features, Trades, Groups
\begin{tikzpicture}[
  hdr/.style={rectangle, draw=isiBlue, fill=isiBlue, text=white,
              minimum width=3.8cm, minimum height=0.55cm,
              align=center, font=\footnotesize\bfseries},
  attr/.style={rectangle, draw=isiBlue!40, fill=white,
               minimum width=3.8cm, minimum height=0.42cm,
               align=left, font=\ttfamily\tiny, inner sep=3pt},
  pk/.style={attr, font=\ttfamily\tiny\bfseries},
  fk/.style={attr, font=\ttfamily\tiny\itshape, text=igsGreen},
  rel/.style={draw=isiBlue!70, thick},
  card/.style={font=\tiny\bfseries\color{isiBlue}}]

  % =================== intern_asset_requests ===================
  \node[hdr] (ir_hdr) at (0,0) {intern\_asset\_requests};
  \node[pk,  below=0pt of ir_hdr] (ir1) {id : UUID (PK)};
  \node[fk,  below=0pt of ir1]   (ir2) {intern\_id : UUID (FK->users)};
  \node[fk,  below=0pt of ir2]   (ir3) {asset\_id : UUID (FK->assets)};
  \node[attr, below=0pt of ir3]  (ir4) {justification : TEXT};
  \node[attr, below=0pt of ir4]  (ir5) {status : ENUM (pending/approved/rejected)};
  \node[attr, below=0pt of ir5]  (ir6) {reviewed\_by : UUID};
  \node[attr, below=0pt of ir6]  (ir7) {voucher\_code : VARCHAR};
  \node[attr, below=0pt of ir7]  (ir8) {reviewed\_at : TIMESTAMPTZ};

  % =================== intern_vouchers ===================
  \node[hdr] (iv_hdr) at (5.2,-1.2) {intern\_vouchers};
  \node[pk,  below=0pt of iv_hdr] (iv1) {id : UUID (PK)};
  \node[fk,  below=0pt of iv1]   (iv2) {request\_id : UUID (FK->requests)};
  \node[attr, below=0pt of iv2]  (iv3) {code : VARCHAR (UNIQUE)};
  \node[fk,  below=0pt of iv3]   (iv4) {created\_by : UUID (FK->users)};
  \node[attr, below=0pt of iv4]  (iv5) {is\_used : BOOLEAN};
  \node[attr, below=0pt of iv5]  (iv6) {used\_at : TIMESTAMPTZ};

  % =================== asset_edit_requests ===================
  \node[hdr] (er_hdr) at (0,-4.8) {asset\_edit\_requests};
  \node[pk,  below=0pt of er_hdr] (er1) {id : UUID (PK)};
  \node[fk,  below=0pt of er1]   (er2) {asset\_id : UUID (FK->assets)};
  \node[fk,  below=0pt of er2]   (er3) {seller\_id : UUID (FK->users)};
  \node[attr, below=0pt of er3]  (er4) {proposed\_changes : JSONB};
  \node[attr, below=0pt of er4]  (er5) {status : ENUM};
  \node[attr, below=0pt of er5]  (er6) {reviewed\_at : TIMESTAMPTZ};

  % =================== trades ===================
  \node[hdr] (tr_hdr) at (5.2,-4.8) {trades};
  \node[pk,  below=0pt of tr_hdr] (tr1) {id : UUID (PK)};
  \node[fk,  below=0pt of tr1]   (tr2) {proposer\_id : UUID (FK->users)};
  \node[fk,  below=0pt of tr2]   (tr3) {receiver\_id : UUID (FK->users)};
  \node[fk,  below=0pt of tr3]   (tr4) {offered\_asset\_id : UUID (FK->assets)};
  \node[fk,  below=0pt of tr4]   (tr5) {requested\_asset\_id : UUID (FK->assets)};
  \node[attr, below=0pt of tr5]  (tr6) {status : ENUM (pending/accepted/declined)};
  \node[attr, below=0pt of tr6]  (tr7) {message : TEXT};

  % =================== groups ===================
  \node[hdr] (gr_hdr) at (10.5,0) {groups};
  \node[pk,  below=0pt of gr_hdr] (gr1) {id : UUID (PK)};
  \node[fk,  below=0pt of gr1]   (gr2) {created\_by : UUID (FK->users)};
  \node[fk,  below=0pt of gr2]   (gr3) {category\_id : INT (FK->categories)};
  \node[attr, below=0pt of gr3]  (gr4) {name : VARCHAR};
  \node[attr, below=0pt of gr4]  (gr5) {description : TEXT};
  \node[attr, below=0pt of gr5]  (gr6) {is\_private : BOOLEAN};
  \node[attr, below=0pt of gr6]  (gr7) {member\_count : INTEGER};

  % =================== group_members ===================
  \node[hdr] (gm_hdr) at (10.5,-4.8) {group\_members};
  \node[fk,  below=0pt of gm_hdr] (gm1) {group\_id : UUID (FK->groups)};
  \node[fk,  below=0pt of gm1]   (gm2) {user\_id : UUID (FK->users)};
  \node[attr, below=0pt of gm2]  (gm3) {role : VARCHAR (member/admin)};
  \node[attr, below=0pt of gm3]  (gm4) {joined\_at : TIMESTAMPTZ};
  \node[attr, below=0pt of gm4]  (gm5) {PRIMARY KEY (group\_id, user\_id)};

  % =================== Relationships ===================
  \draw[rel] (ir_hdr.east) -- node[card,above]{1}++(0.2,0)
                             -- node[card,above]{0..1}(iv_hdr.west);
  \draw[rel] (gr_hdr.south) -- node[card,right]{1}++(0,-0.3)
                              |-node[card,above,near end]{*}(gm_hdr.north);

\end{tikzpicture}
}
\caption{Sprint 4 --- Entity-Relationship Diagram}
\label{fig:sprint4_erd}
\end{figure}

\subsection{Sprint Review --- Deliverables}

\begin{itemize}
  \item Intern asset request workflow with voucher generation and single-use redemption.
  \item Seller analytics dashboard with revenue, download count, and per-asset breakdown.
  \item Asset edit request system allowing sellers to propose post-approval changes.
  \item Peer-to-peer asset trading with accept/decline flow and atomic ownership transfer.
  \item Community groups with member management.
  \item External user seller role request flow with admin approval/rejection.
\end{itemize}

\subsection*{Conclusion}
Sprint 4 delivered the role-specific features that make the RBAC system operationally meaningful for every user segment.
